%=============================================================================!
% 15.6  HOW LONG FOR EACH QUANTITY                                            !
%=============================================================================!
\section{How long for each quantity}
\label{sec:cv-observables}

Each quantity of a run has its own spread and its own memory, so each
needs its own length of run for a given precision. This section measures
both for the liquid, turns them into run lengths, and settles with error
bars the questions that Chapters~12 to 14 left open.

\subsection*{The length of run for a precision}

With $n_{\mathrm{eff}} = t_{\mathrm f}/(2\tau_{\mathrm{int}})$ independent
samples in a run of length $t_{\mathrm f}$ (\cref{sec:cv-correlation}), the
error of the mean is $\sigma_A\sqrt{2\tau_{\mathrm{int}}/t_{\mathrm f}}$.
Asking for an error no larger than a target $\sigma^\ast$ and solving for
$t_{\mathrm f}$,
\begin{equation}
   t_{\mathrm f} \ge \frac{2\tau_{\mathrm{int}}\,\sigma_A^2}{\sigma^{\ast2}} :
   \label{eq:cv-length}
\end{equation}
twice the correlation time for each of the $(\sigma_A/\sigma^\ast)^2$
independent samples needed. Halving the error costs four times
the run. \Cref{fig:cv-observables}(a) checks the $t_{\mathrm f}^{-1/2}$ on the
\SI{2}{\nano\second} run cut into pieces of 50 to \SI{2000}{\pico\second}:
for $U$ the means of the pieces spread by 6.77, 5.14 and 3.36 times the
error of the whole run for pieces of 50, 100 and \SI{200}{\pico\second},
against $\sqrt{40} = 6.32$, $\sqrt{20} = 4.47$ and $\sqrt{10} = 3.16$. The
error bar that each piece finds for itself is lower, 5.38 and 4.19 for the
two shortest,
as \cref{sec:cv-blocks} found, and meets the line from
\SI{200}{\pico\second} on.

\Cref{tab:cv-observables} gives the spread, the correlation time and the
run needed for a stated precision, from \cref{eq:cv-length}, for the
quantities of the liquid. Two things stand out. The pressure has a spread
of \SI{0.0133}{\giga\pascal} per sample about a mean of
\SI{0.085}{\giga\pascal}, a sixth of its value, so with a memory like that
of $U$ and $T$ its mean is slow to fix relative to its size:
\SI{1}{\nano\second} fixes it to 0.7\%, against 0.05\% for $U$. And the
runs at fixed energy need a 15th to a 38th of the time needed under CSVR
for the same precision of $U$, $T$ and $P$, because their correlation
times are 8 to 14 times shorter (\cref{sec:cv-correlation}) and their
spreads smaller. A run at fixed energy samples the microcanonical
ensemble, at the temperature its energy fixes: here \SI{133.70}{\kelvin},
\SI{1.1}{\kelvin} below the CSVR run, because the energy it was given is
the mean of \SI{40}{\pico\second} of CSVR, which, with the
\SI{1.71}{\pico\second} memory of the energy, is uncertain by about
\SI{1.6}{\kelvin} of temperature.
\Cref{sec:th-compare} recommended preparing a run under a thermostat and
gathering its dynamical quantities at fixed energy; for averages too,
fixed energy reaches a precision sooner, at the temperature its energy
sets.

\begin{table}[htb]
\centering
\small
\caption{How long each quantity needs, for liquid argon at
\SI{135}{\kelvin}: the spread of one sample, the integrated correlation
time under CSVR with $\tau_{\mathrm T} = \SI{1}{\pico\second}$ and at fixed
energy, the error of the mean after \SI{1}{\nano\second} under CSVR, and
the run needed for the target error $\sigma^\ast$, from
\cref{eq:cv-length}. The
volume is from Chapter~14's runs under stochastic cell rescaling at
\SI{0.1}{\giga\pascal}; $D$ from blocks of \SI{100}{\pico\second}.}
\label{tab:cv-observables}
\begin{tabular}{@{}lccccccc@{}}
\toprule
& spread & \multicolumn{2}{c}{$\tau_{\mathrm{int}}$ / ps} & error of
& target & \multicolumn{2}{c}{run / ns}\\
& CSVR / fixed $E$ & CSVR & fixed $E$ & \SI{1}{\nano\second} &
$\sigma^\ast$ & CSVR & fixed $E$\\
\midrule
$U/N$ & \SI{0.677}{\milli\electronvolt} / 0.532 & 0.73 & 0.077 &
\SI{0.026}{\milli\electronvolt} & \SI{0.01}{\milli\electronvolt} & 6.7 & 0.44\\
$T$ & \SI{6.93}{\kelvin} / 4.13 & 1.05 & 0.077 & \SI{0.32}{\kelvin} &
\SI{0.1}{\kelvin} & 10 & 0.26\\
$P$ & \SI{0.0133}{\giga\pascal} / 0.0096 & 0.94 & 0.123 &
\SI{0.00058}{\giga\pascal} & \SI{0.001}{\giga\pascal} & 0.33 & 0.023\\
$\mathsf{P}_{xy}$ & \SI{0.0121}{\giga\pascal} & 0.18 & 0.18 & & & &\\
$V$ & \SI{185}{\angstrom\cubed} & 2.5 & & \SI{13}{\angstrom\cubed} &
\SI{10}{\angstrom\cubed} & 1.7 &\\
$D$ & & & & 1.0\% & 1\% & 1.0 &\\
\bottomrule
\end{tabular}
\end{table}

\subsection*{Diffusion under a thermostat}

Chapter~13 found the diffusion coefficient under CSVR, Nosé-Hoover chains
and Berendsen coupling 1.01 to 1.07 times its value at fixed energy, in
runs 0.4 to \SI{1.8}{\kelvin} warmer, and left open whether the difference
was real (\cref{sec:th-compare}). Cutting each run of this chapter into
blocks of \SI{100}{\pico\second}, long beside the \SI{20}{\pico\second}
over which $D$ is fitted, gives independent estimates whose spread gives
the error. The \SI{2}{\nano\second} run under CSVR has $D = (3.877 \pm
0.028)\times10^{-4}$\,\si{\angstrom\squared\per\femto\second} at $(134.81
\pm 0.22)$\,\si{\kelvin}; the \SI{2}{\nano\second} run at fixed energy has
$(3.834 \pm 0.034)\times10^{-4}$ at \SI{133.70}{\kelvin}. The two are at
different temperatures, and runs under CSVR at 130 and \SI{140}{\kelvin}
give the slope $\dd D/\dd T = (0.0294 \pm 0.0071)\times10^{-4}$\,\si{\angstrom
\squared\per\femto\second\per\kelvin}
(\cref{fig:cv-observables}(b)). Moved along it to \SI{134.81}{\kelvin},
the run at fixed energy gives $(3.867 \pm 0.035)\times10^{-4}$, and the
ratio of the two is $1.002 \pm 0.011$. Berendsen coupling with
$\tau_{\mathrm T} = \SI{0.1}{\pico\second}$ at 130 and \SI{140}{\kelvin}
lies 1.3 and 0.3 errors from the line through the CSVR runs. To about 2\%,
CSVR with $\tau_{\mathrm T} = \SI{1}{\pico\second}$ and Berendsen coupling
with $\tau_{\mathrm T} = \SI{0.1}{\pico\second}$ leave diffusion alone.
Chapter~13 found $1.05 \pm 0.01$ and $1.04 \pm 0.01$ for these two from
three runs of \SI{100}{\pico\second} each. A block of
\SI{100}{\pico\second} here fixes $D$ to 3.2\%, so a mean of three such
runs is uncertain by $3.2/\sqrt3 = 1.8\%$ and a ratio of two such means
by $1.8\times\sqrt2 = 2.6\%$: the two values lie within about two such
errors of 1 once their warmth is removed, and the $\pm0.01$ came from the
spread of only three runs, itself uncertain by half
(\cref{ex:cv-error-error}). The other couplings of Chapter~13 were not run
again.

\begin{figure}[htb]
\centering
\includegraphics{ch15_convergence/observables}
\caption{(a) The error of the mean of $U$, $T$ and $P$ against the length
of run: the \SI{2}{\nano\second} CSVR run cut into pieces, the error bar
each piece finds for itself averaged over the pieces (filled) and the
spread of the pieces' means (open), each divided by the error of the whole
run, against the inverse square root of the length (dashed). (b) The
diffusion coefficient against
the temperature, with errors from blocks of \SI{100}{\pico\second}: CSVR
(teal), Berendsen coupling with $\tau_{\mathrm T} = \SI{0.1}{\pico\second}$
(ochre) and fixed energy (grey), with the line through the CSVR runs at
130 and \SI{140}{\kelvin} (dashed).}
\label{fig:cv-observables}
\end{figure}

\subsection*{Questions left open}

Chapter~12 found the light atoms of a mixture at \SI{126}{\kelvin} and
the heavy at \SI{124}{\kelvin} over the last \SI{10}{\pico\second} of a run,
and asked whether the gap was real (\cref{sec:sm-equipartition}). Taking
the difference of the two temperatures at each recorded frame as the
series, its mean is $(1.30 \pm 2.14)$\,\si{\kelvin}, with $g = 31.4$: 0.6
errors from zero, and equipartition between the species holds within what
the run can tell. Chapter~12 also found single runs of \SI{50}{\pico\second}
at fixed energy to give heat capacities scattered by up to 7.5\% about the
slope of $E$ against $T$, $2.376\,N\kB$ (\cref{sec:sm-kinetic}); the five
runs of 256 atoms lie within 5.3\% of it, and the 7.5\% is a run of 864
atoms. A heat capacity comes from a variance, whose error the rule for a
mean does not give; it is found by resampling, which \cref{sec:cv-ensemble}
sets out. It gives the five runs of 256 atoms errors of 2.9 to 3.9\%,
$(2.292 \pm 0.067)\,N\kB$ near \SI{130}{\kelvin} among them, and all five
lie within 1.7 of their errors of the slope: their scatter was the noise
of the runs.
Chapter~14 quoted the pressure of the liquid as \SI{0.0877}{\giga\pascal}
(\cref{sec:pr-periodic}), from 50 frames \SI{2}{\pico\second} apart,
farther apart than $2\tau_{\mathrm{int}} = \SI{1.9}{\pico\second}$, so
nearly independent: its error is
$0.0133/\sqrt{50} = \SI{0.0019}{\giga\pascal}$. The \SI{2}{\nano\second}
run gives $(0.0849 \pm 0.0004)$\,\si{\giga\pascal}, 1.4 combined errors
away.

\begin{keyidea}
A run of length $t_{\mathrm f}$ fixes a mean to
$\sigma_A\sqrt{2\tau_{\mathrm{int}}/t_{\mathrm f}}$, so a target error
$\sigma^\ast$ needs $t_{\mathrm f} \ge 2\tau_{\mathrm{int}}
\sigma_A^2/\sigma^{\ast2}$. Each quantity has its own $\sigma_A$ and
$\tau_{\mathrm{int}}$, and a thermostat lengthens the memory of everything
tied to the energy. With error bars, the diffusion coefficient is the same
under CSVR with $\tau_{\mathrm T} = \SI{1}{\pico\second}$, Berendsen
coupling with $\tau_{\mathrm T} = \SI{0.1}{\pico\second}$ and fixed
energy to about 2\%.
\end{keyidea}

\notebook{15}{find the run each quantity of the liquid needs}

\begin{selfcheck}
A quantity has $\sigma_A = 2$ and $\tau_{\mathrm{int}} = \SI{5}{\pico
\second}$. How long a run fixes its mean to $\pm0.1$?
\end{selfcheck}
